%!TEX root = ../main.tex

\chapter{Diseño del software del PFD}
\label{ch:software}

El \cref{ch:firmware} ha descrito cómo el nodo receptor entrega al PC el estado
de vuelo decodificado mediante la frase \texttt{\$TEL}. Este capítulo
describe el software de la estación de tierra, una aplicación que recibe los datos y representa los parámetros en un PFD con disposición Basic-T (RF09, RN01). Se describen la disposición y simbología
del PFD, la recepción y el suavizado de los datos, los instrumentos y, por
último, la vista de mapa de navegación.


\section{Entorno y organización}
\label{sec:sw_entorno}

El PFD se implementa en Processing,\footnote{Processing, entorno de
\textit{software} libre para programación gráfica: \texttt{processing.org}.} un
entorno de programación gráfico de código abierto basado en Java. Se elige por su gestión directa del dibujo en dos dimensiones,
su portabilidad y su baja barrera de entrada, adecuada para una herramienta
docente. El render utiliza el motor P2D a \SI{60}{FPS}. Se escoge este motor, ya que se observan menos defectos en pantalla comparado con el motor por defecto.

El código se organiza en un sketch principal (\texttt{FlightDisplay.pde}), que
gestiona la recepción por puerto serie, el bucle de dibujo y el escalado de la
ventana, y un archivo por cada instrumento: actitud, altímetro, VSI, velocidad
respecto al suelo, rumbo y proximidad. Un archivo adicional
(\texttt{GPSMap.pde}) añade la vista de mapa.


\section{Disposición y simbología del PFD}
\label{sec:sw_basict}

La disposición sigue el patrón Basic-T de los displays electrónicos de vuelo
\cite{ac25_11b}, equivalente al del \textit{glass cockpit} de aviación general.
El indicador de actitud ocupa la posición central y de mayor tamaño, a la
izquierda se sitúa la velocidad respecto al suelo y a la derecha el altímetro y
el VSI. En la parte inferior del conjunto, el indicador de rumbo y en la esquina inferior, el indicador
de proximidad. En la \cref{fig:basict_layout} se representa de forma general cómo se organizan los indicadores.

\begin{figure}[H]
    \centering
    \begin{tikzpicture}[font=\footnotesize,
        inst/.style={draw, rounded corners=2pt, align=center, fill=black!5}]
        \node[inst, minimum width=34mm, minimum height=34mm] (ai) at (0,0)
            {Indicador\\de actitud\\(horizonte\\artificial)};
        \node[inst, minimum width=12mm, minimum height=34mm, left=6mm of ai] (gs)
            {GS\\(vel.\\suelo)};
        \node[inst, minimum width=12mm, minimum height=34mm, right=6mm of ai] (alt)
            {Altímetro};
        \node[inst, minimum width=8mm,  minimum height=34mm, right=2mm of alt] (vsi)
            {VSI};
        \node[inst, minimum width=52mm, minimum height=11mm, below=7mm of ai] (hdg)
            {Indicador de rumbo};
        \node[inst, minimum width=22mm, minimum height=11mm, right=6mm of hdg] (prox)
            {Proximidad};
    \end{tikzpicture}
    \caption{Disposición Basic-T de los instrumentos del PFD (RN01): actitud
             central, velocidad de tierra a la izquierda, altímetro y VSI a la
             derecha, rumbo debajo y proximidad en la esquina.}
    \label{fig:basict_layout}
\end{figure}

La simbología y el código de color implementados siguen las guías de displays electrónicos y
de factores humanos \cite{sae_as8034c, hfstd001b} y la jerarquía de color para
alertas de la normativa \cite{cfr25_1321, cfr25_1322} (RN02, RN03). Estos colores están resumidos en
la \cref{tab:sw_colores}.

\begin{table}[H]
    \centering
    \caption{Código de color de la simbología del PFD}
    \label{tab:sw_colores}
    \begin{tabularx}{\linewidth}{l l X}
        \toprule
        \textbf{Elemento} & \textbf{Color} & \textbf{Uso} \\
        \midrule
        Cielo / tierra          & Azul / marrón & Horizonte artificial. \\
        Simbología fija         & Blanco        & Escalas, marcas y símbolo del avión. \\
        Lecturas numéricas      & Verde         & Valores actuales. \\
        Referencias ajustables  & Cian          & \textit{Bug} de rumbo y ajuste Kollsman. \\
        Precaución              & Ámbar         & Zona intermedia (proximidad, VSI). \\
        Peligro                 & Rojo          & Enlace perdido y proximidad crítica. \\
        \bottomrule
    \end{tabularx}
\end{table}


\section{Recepción y flujo de datos}
\label{sec:sw_recepcion}

La frase \texttt{\$TEL} se recibe por el puerto serie y se procesa línea a
línea. Cada campo actualiza un \emph{valor objetivo}, no el valor mostrado:
cabeceo, alabeo, rumbo, altitud, velocidad vertical, posición, satélites,
indicadores de estado y la distancia de proximidad. El bucle de dibujo, a
\SI{60}{FPS}, interpola el valor mostrado hacia el objetivo
(\cref{sec:sw_suavizado}). Si transcurren más de \SI{2}{\second} sin recibir
paquetes, el enlace se marca como perdido y el PFD lo señala. Las tramas
corruptas (por interferencias o bytes perdidos) se descartan sin interrumpir la
representación.


\section{Suavizado de la representación}
\label{sec:sw_suavizado}

El paquete llega cada \SI{20}{\hertz}, pero el PFD se redibuja a \SI{60}{FPS}. Para
que el movimiento sea continuo, cada parámetro se interpola linealmente hacia su
objetivo en cada fotograma (filtro EMA):
\begin{equation}
    x[k] = x[k-1] + \alpha\,(x_{\text{obj}} - x[k-1]), \qquad \alpha = 0{,}4
\end{equation}
 Para algunos parametros también se implementa un valor mínimo antes de mostrar cambios, para evitar ver \textit{jitter} en pantalla. Este suavizado enmascara la cadencia del enlace sin añadir un retardo perceptible.


\section{Instrumentos}
\label{sec:sw_instrumentos}

\paragraph{Indicador de actitud.}
El horizonte artificial se dibuja como dos rectángulos (cielo y tierra, con
gradiente) rotados según el alabeo y desplazados según el cabeceo, sobre los que
se superponen la escala de cabeceo y el arco de alabeo. Una máscara circular (un
rectángulo con un agujero circular recortado) deja visible solo el interior del
instrumento. La referencia de alabeo es fija al instrumento (estilo
\textit{sky pointer}).

\paragraph{Altímetro.}
Cinta vertical de altitud MSL con una ventana de Kollsman que muestra el ajuste
barométrico, modificable por teclado. Este instrumento se modifica respecto al estándar, ya que en este caso el desplazamiento vertical es muy inferior a lo normal. Se hace el rango visible inferior para que el movimiento sea apreciable.

\paragraph{Velocidad vertical (VSI).}
Cinta con una barra de tendencia cuyo color varía con la magnitud (verde, ámbar,
naranja). Igual que en el altímetro, la escala también se reduce para uso en banco.

\paragraph{Velocidad respecto al suelo.}
Cinta vertical alimentada por la velocidad derivada del GPS, convertida a nudos. Este instrumento casi no tiene utilidad en esta aplicación, pero se implementa para cumplir con el estándar de PFD.

\paragraph{Indicador de rumbo.}
Cinta horizontal con marcas cardinales y un \textit{bug} de rumbo cian ajustable
con las flechas del teclado. Las marcas mayores se etiquetan según la convención
aeronáutica (090$^{\circ}$ se muestra como ``9'').

\paragraph{Indicador de proximidad.}
Muestra la distancia al suelo con un código de color por umbrales (rojo, ámbar y
verde, además de los estados de seguridad y de ausencia de datos). De forma
coherente con la jerarquía de la \cref{sec:fw_proximidad} (RF06), la fuente
activa es el ToF. El código que pertenece al radar también existe, pero se descarta el uso de sus datos debido a la poca fiabilidad que ofrece. El conjunto final se muestra en la
\cref{fig:pfd_screenshot}.

\begin{figure}[H]
    \centering
    \includegraphics[width=\linewidth]{Figures/pfd.png}
    \caption{PFD en disposición Basic-T con telemetría en vivo.}
    \label{fig:pfd_screenshot}
\end{figure}


\section{Vista de mapa de navegación}
\label{sec:sw_mapa}

Como funcionalidad adicional al Basic-T, el PFD incorpora una vista de mapa de
navegación, conmutable por teclado, que aprovecha la posición GPS (RF05). El
mapa proyecta la posición en Web Mercator y emplea el esquema de teselas de OpenStreetMap.\footnote{Convención de teselas
\textit{slippy map} de OpenStreetMap} Cada
tesela se identifica por su nivel de zoom y sus índices, y se sitúa en pantalla
mediante la lógica siguiente:
\begin{equation}
    x = \frac{\lambda+180}{360}\,256\cdot 2^{z}, \qquad
    y = \frac{1}{2}\left[1 - \frac{\ln\!\left(\tan\varphi + \sec\varphi\right)}{\pi}\right]256\cdot 2^{z}
\end{equation}
con $\lambda$ la longitud, $\varphi$ la latitud y $z$ el nivel de zoom.

Las teselas son topográficas del ICGC,\footnote{Servei de Mapes Base de l'ICGC,
capa topográfica; licencia CC-BY (Cataluña) / ODbL (resto del mundo):
\texttt{icgc.cat}.} descargadas previamente a disco con un script en Python para que el mapa funcione sin
conexión. Las teselas ausentes, no descargadas, se dibujan como celdas grises. Sobre el mapa se trazan la traza GPS (normalmente no es visible, ya que no hay desplazamiento suficiente de la placa), el símbolo del avión orientado al rumbo actual, una
rosa de los vientos, una barra de escala calculada según la resolución real en
metros por píxel de Web Mercator, y por último una barra con información general de los datos del GNSS. La disposición de esta vista sigue las recomendaciones de \cite{hfstd001b} para
displays de mapa. El resultado se muestra en la \cref{fig:gps_map}.

\begin{figure}[H]
    \centering
    \includegraphics[width=\linewidth]{Figures/map.png}
    \caption{Vista de mapa de navegación con teselas topográficas del ICGC, traza
             GPS y símbolo del avión orientado al rumbo.}
    \label{fig:gps_map}
\end{figure}
